\section{Open Weights Go Full-Lifecycle}
\sectionkicker{OPEN EFFICIENT}
\begin{wideflow}
{\Large\bfseries 開いた重みは作りごと競う}\par
\smallskip
{\color{SurveyMuted}全物公開の連なりと境界付き重みを、許しと日付の精度ごと示す。}\par
\end{wideflow}

9月3日にIFMが出したK2 Horizonは、375B-A23B、36B-A4B、32B、7B、3.7B、0.9Bの6連なりで、Hugging Faceの収集置き場とGitHubの符号、Weights \& Biasesに物が置かれる\cite{k2horizon}。模型と符号はApache 2.0で出るが、資料群は適用される条件（例ODC-BY）に従い、作り方や混ぜ方の開示は再配布不能の範囲で添えられる。Apache 2.0を資料群に広げて書かない。6規模とも量子化に対応し、初日からvLLM、SGLang、Ollamaが動くとされる。提供側の順位主張は、0.9B・3.7B・7Bが各規模で最上、36B-A4BがMoVAの疎注意設計で約40億の動作素子により密32B級に迫る、32Bと375B-A23Bが各級の上位という内容である。いずれも提供側の数値であり、独立の再現を待つ。MoVAの効率解釈も提供側の設計主張であり、給仕の独立検証はない。開放の内訳は、学習資料や資料調合法、学習符号、設定、中間検査点、細粒度の記録、評価結果、最終重みが規模ごとに示され、事前学習は模型あたり約20T token、事後学習は1億超の仕事とされる。TerminalBench 2.1の監査訂正（旗付き24試行を除き70.2\%から66.9\%へ）が開示された点は記録の範囲で押さえる。

\begin{claimboundary}[許しの境界]
Apache 2.0は模型と符号の話であり、資料群に広げない。順位と効率の主張は提供側の報告であり、独立の再現を待つ。
\end{claimboundary}

8月28日に置かれたzai-orgのGLM-5.3は全重み（753B素子、BF16・F8\_E4M3・F32のsafetensors）を開き、Transformers、vLLM、SGLang、KTransformers、Unsloth、Dockerの給仕経路と、reasoning\_effortのlow・high・max（既定max）、論文arXiv 2602.15763への言及が札に示される\cite{glm53weights}。GLM-5.2からの伸びとして、自社Z.ai Code Benchでの50\%増、Terminal Bench 3.0とAgents' Last Examの公開最上、CyberGym 84.5と悪用連鎖の倍増、Terminal Bench 2.1の88.2、AutomationBenchの48.2、GDPval-AA v2の1769といった表値はすべて札の報告であり、外部の再現はない。札のLicense欄はglm-5.3（独自）と読み、Flash系に付くMITの許しを全重みに借りてこない。全文の許し文は消費していない。取得のmarkdownに明示の公表日がなく、8月28日の協定時単位は未解決のため、窓配置は中確度の境界扱いとし、時単位の断定はしない。落としや好みの数は基盤の遠隔測定であり、能力の証拠にはしない。

\begin{technicalnote}[日付と許しの注意]
8月28日の協定時単位は未解決であり、窓配置は中確度にとどめる。Flash系のMITを全重みに持ち込まない\cite{glm53weights}。
\end{technicalnote}
